\chapter{Implementazione Filtri FIR e IIR}

\begin{figure}[h]
    \centering
    \includegraphics[scale = 0.8]{Signal_flow_graph_for_a_digital_IIR_Biquad.png}
\end{figure}  

\newpage 

\section{Ripasso degli argomenti precedenti}

\footnote{ 
    Slide del prof | 11 - Implementazione Filtri FIR e IIR 1 | pag 2 - 4\\
    Appunti | 2025-11-25 | pag 2 \\ 
    Latex | 2025-12-16 Mattino | pag 280   
} 

Prima di descrivere degli SFG (Signal Flow Graph) è necessario svolgere un piccolo ripasso degli argomenti precedenti. \newline 

Come studiato in \ref{sec: Sistemi LTI descritti da equazioni alle differenze} \nameref{sec: Sistemi LTI descritti da equazioni alle differenze} a pag \pageref{sec: Sistemi LTI descritti da equazioni alle differenze} 
ed in \ref{sec: Funzioni di Rete} \nameref{sec: Funzioni di Rete} a pag \pageref{sec: Funzioni di Rete}, 
una funzione $H \left(z\right)$ razionale corrisponde ad un'equazione lineare alle differenze finita a coefficienti costanti. \newline 

Per questa classe di funzioni è possibile individuae 3 blocchi funzionali: 

\begin{figure}[h]
    \centering
    \includegraphics[scale = 0.4]{blocchi funzionali.png}
\end{figure}  

Come studiato ad Elettrotecnica per i circuiti reali, anche qui a DSP, data la stessa equazione, sono possibili diverse realizzazioni circuitali. \newline 

Le caratteristiche da ricercare sono: 

\begin{itemize}
    \item minima complessità perchè, per un calcolatore, meno conti si fanno e meglio è  
    \item minima sensibilità al troncamento, cioè circuiti più robusti alla quantizzazione 
\end{itemize}

\newpage  

\section{Signal Flow Graph}

\footnote{ 
    Slide del prof | 11 - Implementazione Filtri FIR e IIR 1 | pag 5 - 6\\
    Appunti | 2025-11-25 | pag 2  
} 

Come studiato ad Elettrotecnica, useremo i grafi per rappresentare le equazioni, quindi i circuiti. \newline 

Ma, a differenza di Elettrotecnica dove i grafi sono bidirezionali, cioè lo stesso ramo può essere percorso in tutti i due i versi, 
negli SFG (Signal Flow Graph) i grafi sono unidirezionali. \newline 

Non solo: ad Elettrotecnica ogni ramo conteneva due informazioni, cioè tensione e corrente, invece qui ogni ramo contiene solo un'informazione. \newline 




